% ============================================================
%  02_introduction.tex  –  Chapter 1: Introduction
%  Template: chapter title 14 pt Bold Uppercase Center
%            section title 12 pt Bold Uppercase Left
%            text 12 pt Arial justified, 1.5 cm first-line indent
% ============================================================

\chapter{Introduction}
\label{ch:introduction}

\section{Context}
\label{sec:context}

Air quality always had a huge impact on human health. Although it has gained huge attention as
time passed, another similar aspect, which can be considered just as important, is indoor air
quality. People spend most of their time indoors, so having the space properly monitored is of
great importance. Especially when we consider the fact that indoor spaces have limited
ventilation and usually have multiple people in them, which results in poor air quality.

Usual activities, such as cleaning, cooking or using heating appliances can lead to the release
of gases and harmful substances. Beside these, ambient noise and temperature and humidity level
can pose a threat to indoor comfort. Commercial air monitoring solutions exist for these
problems, although they provide limited accessibility to the user. They usually offer limited
possibilities to the user to access the raw data or long-term historical data. Same with
applying custom predictive analytics. This shows us that transparent, customizable products
which provide real-time data and its analysis are more than welcome and wished for.

\section{Motivation}
\label{sec:motivation}

The main motivation for this project is the desire to create an indoor air monitoring system
that connects raw hardware data collection and software. While sensors provide initial metrics,
the true values are generated when this data is computed, historically saved and made easily
accessible to the user. This project presents an attempt at unifying different technologies:
hardware using Raspberry Pi and sensors, web applications, databases and machine learning.

From a health and sociological perspective, monitoring these factors on an ongoing basis can be
particularly relevant in preventing problems related to respiratory system and ensuring healthy
sleeping habits.

The objective of this project is to motivate taking action through the presented data by
presenting the tools for taking the initiative in one's environment and taking control. This
project intends to provide the user with the means necessary to collect and transfer his/her
environmental information without relying on proprietary technology.

\section{Personal Contributions}
\label{sec:contributions}

The main contribution of this thesis is that of a system that integrates all stages of air
quality detection: sensor measurement, cloud database, user website and mobile interface for
data visualization and machine learning prediction, all in one project. In comparison,
commercial products like Awair Element and uRADMonitor A3 provide a more advanced and
user-friendly interface (either through browser or mobile interface), but the data remains
proprietary. Sensors cannot be added or replaced either. Literature research into
research-level multi-parameter air quality monitoring systems, such as those created by
Ganesan et al.\ and Kim et al., shows the potential of the same kind of sensors, but does not
include a full-stack solution with user authentication and predictions.

\airnest{} combines two layers of real-time alerts, providing both the immediate
threshold-based alerts generated by the system at the hardware level on a Raspberry Pi and the
7-day forecast based on machine learning available on the public website. Predictions are made
based on a hybrid model using three data sources: historical Open-Meteo weather data, historical
EEA government monitoring station readings for Timișoara and live readings from the station.
This is something not implemented by any of the systems reviewed.
